% Results and Analysis. Requires the main table (Table~\ref{tab:main}), the
% scaling table (Table~\ref{tab:scaling}), and the appendix full matrices
% (tab:single-node-235b, tab:full-matrix-35b, tab:full-matrix-b200).

\section{Results and Analysis}
\label{sec:results}

Table~\ref{tab:main} reports single-node SLO-qualified goodput for both models on
both vendors; Table~\ref{tab:scaling} reports multi-node scaling; the full
per-load-point matrices are in the appendix. Read together, the runs tell a
consistent story that we organize as five insights. The recurring theme is that
\emph{long-context serving capacity is governed by where the prefill and
tail-latency bottlenecks fall and by the maturity of the serving stack, far more
than by raw accelerator FLOPS}---and that ranking on SLO-qualified goodput, rather
than raw throughput, is what makes these effects visible.

\subsection{Sparsity buys concurrency, not a faster request}
\label{subsec:insight-concurrency}
The headline gap between the tiers is enormous---at LC-8K the 35B-A3B reaches
$1451$\,tok/s (SGLang) against $31$ for the 235B-A22B, a $\sim$47$\times$ ratio,
mirrored by vLLM ($952$ vs.\ $22$)---but it is \emph{not} a faster request. At
concurrency~1 the per-stream rates differ by only $\sim$2$\times$ ($30.7$ vs.\
$14.7$\,tok/s, vLLM). The entire remaining gap is \textbf{SLO-qualified
concurrency}: the 22\,B-active 235B violates the $120$\,ms TPOT bound by
concurrency~4 (TPOT$_{p95}{=}153$\,ms), capping at $\text{C}^\star{=}2$, whereas
the 3\,B-active 35B holds TPOT$_{p95}\!\leq\!67$\,ms to concurrency~64---a
$32\times$ larger batch inside the same tail-latency envelope. The practical
reading is that MoE sparsity is a \emph{batchability} lever: it does little for
single-request latency but makes batched decode far cheaper, which under a
tail-latency SLO converts directly into order-of-magnitude interactive capacity
per node (at the cost of headline quality, \S\ref{subsec:insight-goodput}). This
also explains the energy column: because mean power varies little across configs
($340$--$534$\,W), tokens/joule tracks goodput, so the most batchable
configuration (SGLang) is also the most energy-efficient ($0.419$ vs.\ $0.100$
tok/J on the 35B LC-Cache).

\subsection{Long-context serving is prefill-bound, and the wins are in scheduling and reuse}
\label{subsec:insight-prefill}
Three otherwise-unrelated results share one cause---prefill, not decode, is the
long-context bottleneck. (i)~\textbf{The ``128K wall'' is a scheduling property.}
SGLang is the only backend to satisfy the heavy SLO at LC-128K, on both the 235B
($0.9$\,tok/s) and the 35B ($257$ at C$^\star{=}16$); vLLM fails LC-128K on both. The cause is
prefill latency: vLLM's 128K TTFT$_{p95}\!\approx\!20$\,s exceeds the $10$\,s
bound, while SGLang's chunked prefill brings it to $\approx7$\,s. (ii)~\textbf{The
value of a prefix cache scales with how prefill-bound the model is.} Disabling the
prefix cache (cold vs.\ warmed) collapses LC-Cache goodput $\sim$15$\times$ on the
235B ($\approx47\!\to\!3$) but only $\sim$1.1$\times$ on the 35B ($34\!\to\!30$):
the 3\,B-active model re-prefills a 32K context so cheaply that caching saves
little, while the 22\,B-active model is prefill-bound and benefits enormously.
(iii)~\textbf{Faster prefill silicon clears the wall:} on B200, vLLM \emph{does}
qualify LC-128K ($\approx1.2$, Table~\ref{tab:main}) where it failed on MI355X---a
prefill-FLOPS effect. The unifying lesson is that long-context throughput is won by
attacking prefill (chunked-prefill scheduling, prefix reuse, prefill FLOPS), and
that graph capture---which optimizes only the decode launch path---cannot help
here; indeed, enabling HIP graphs on gfx950 \emph{regressed} throughput and failed
to complete the long-context ladders, so all AMD numbers use eager mode.

\subsection{The serving stack dominates the hardware}
\label{subsec:insight-stack}
Holding the model and GPU fixed, the software choices move goodput more than the
silicon does. \textbf{Backend:} SGLang beats vLLM on every workload of both models.
\textbf{Precision:} \emph{dynamic} FP8 \emph{lowers} goodput while preserving
quality (235B: $16.8$ vs.\ $22.2$ at LC-8K; 35B at TP1: $149$ vs.\ $195$) and has
the worst tokens/joule in the table---its on-the-fly dequantization is not
amortized---whereas \textbf{offline MXFP4} (AMD Quark) on the 235B \emph{matches}
BF16 goodput at $\sim$9\% lower power, the best tokens/joule among the vLLM
precisions ($0.110$ vs.\ $0.098$). So block-FP4 from a pre-quantized checkpoint, not
dynamic FP8, is the FP4 efficiency lever. \textbf{Routing:} for multi-node
scale-out the routing \emph{policy}, not the GPU count, governs the result
(Table~\ref{tab:scaling}). On the single-shared-prefix LC-Cache workload,
cache-aware routing keys on the prefix and pins essentially all traffic to one
replica---throughput stays at the single-replica $68.3$\,tok/s for every $N$, i.e.\
$\text{Efficiency}(N){=}1/N$ ($0.25$ at $N{=}4$, $0.13$ at $N{=}8$)---while
round-robin spreads load and scales near-linearly ($1.00$ at $N{=}2$, $0.75$ at
$N{=}4$), at the cost of redundant per-replica prefix prefill. This is a concrete
cache-locality vs.\ load-balance tradeoff that no single policy resolves, arguing
for routers that shard distinct prefixes while replicating hot ones. Finally,
\textbf{stack maturity} is itself a first-order variable: TGI's ROCm images did not
serve gfx950, the vLLM fused-MoE MXFP4 loader could not read the 35B checkpoint, and
the 235B's dual-chunk 1M context did not load on ROCm---each a coverage gap that
bounds what the hardware can be asked to do today.

\subsection{SLO-qualified goodput exposes wins that raw throughput hides}
\label{subsec:insight-goodput}
Two results justify ranking on quality-gated, SLO-qualified goodput rather than raw
tokens/s. First, the 35B is fast but \emph{wrong} on cached RAG: it is
latency-SLO-qualified on LC-Cache yet \textbf{fails the needle-recall guardrail}
(QualityRatio $0.2$--$0.6$ vs.\ $1.0$ for the 235B), uniformly on single-node and
multi-node and independent of concurrency---the 3\,B-active model serves the 32K
shared-prefix queries quickly but answers the embedded needle incorrectly most of
the time. Raw throughput would reward this; quality-gated goodput does not, so we
report the 35B LC-Cache figures as valid-only (marked $\ddagger$ in
Table~\ref{tab:main}). Second, the same gate clarifies scaling: the round-robin
\emph{raw} throughput scales cleanly ($136.5\!\to\!204.8$ tok/s from $N{=}2$ to
$4$), but the quality factor is what we must report alongside it. The broader point
is that for the small MoE the speed-vs-quality tradeoff is not hypothetical---it is
the difference between a leaderboard-topping number and a usable system---and the
metric design is what surfaces it. The two retrieval-shape workloads
(Table~\ref{tab:rag}) make the same point on \emph{realistic} RAG. On RAG-TopK the
35B's throughput batches well ($17\!\to\!83$\,tok/s raw to C$^\star{=}8$) but it
answers only $\sim$20\% correctly (quality $0.17$--$0.21$ across both backends),
retrieving the wrong magic number from among the distractor passages, whereas the
235B is correct (quality $1.00$) at a far lower $3.6$--$5.4$\,tok/s (prefill-bound at
32K). Crucially, even the 35B's \emph{quality-gated} goodput ($17$) exceeds the
235B's ($3.6$--$5.4$): neither raw \emph{nor} quality-gated goodput rejects a model
that is wrong $\sim$80\% of the time---only the separately reported quality score
does, which is the strongest argument for treating quality as a first-class metric.
RAG-Report, by contrast, is answered correctly by both tiers but is decode-bound,
and exposes a $\sim$50$\times$ \emph{backend} gap in long-output batching: 35B SGLang
reaches $501$\,tok/s at C$^\star{=}32$ while 35B vLLM stalls at $10$ (C$^\star{=}1$).
The two workloads thus separate a \emph{retrieval-accuracy} failure from a
\emph{decode-capacity} limit---distinct dimensions that raw tokens/s conflates.

\paragraph{External validity on public real-text QA.}
Because the synthetic corpus buys control at the cost of realism, we repeat the
retrieval comparison on two public LongBench workloads---real Wikipedia
multi-document QA (RAG-RealQA, HotpotQA) and scientific-paper QA (RAG-SciQA,
QASPER), graded by token-F1 (Table~\ref{tab:realqa})---and they sharpen, rather than
contradict, the quality argument. First, \emph{the synthetic needle test overstates
retrieval quality}: the 235B scores F1$=1.0$ on synthetic RAG-TopK but only $0.58$
on real Wikipedia QA and $0.25$ on scientific QA, so a leaderboard built only on
synthetic exact-match would report a quality ceiling that real traffic never
reaches---direct evidence that quality must be measured on real corpora, not just
controlled ones. Second, \emph{the quality ordering between tiers is
workload-dependent}: on real multi-hop QA the smaller 35B is actually
competitive-to-better (F1 $0.73$--$0.75$ vs.\ $0.57$--$0.58$ once its reasoning mode
is disabled so the short-answer budget is not spent on \texttt{<think>} tokens),
inverting the large synthetic-RAG-TopK gap, while scientific QA is hard for both
tiers ($0.16$--$0.28$). The synthetic ``fast-but-wrong'' separation is therefore
specific to \emph{adversarial same-format needle} retrieval and does not
generalize to all real RAG; the throughput story is unchanged (the more batchable
35B sustains higher goodput, the prefill-bound 235B caps at C$^\star{=}1$). The
practical lesson reinforces the metric design: quality is first-class \emph{and}
corpus-dependent, so a serving benchmark must pair controlled synthetic shapes with
public real-text workloads to avoid both contamination and over-optimism.

\subsection{The cross-vendor verdict is workload-dependent}
\label{subsec:insight-crossvendor}
Running the 235B on NVIDIA B200 with the identical harness
(Table~\ref{tab:main}) yields no single ordering. B200 raises the short/medium
workloads (vLLM LC-8K $22.2\!\to\!36.4$, LC-32K $3.6\!\to\!7.2$) and, as above,
uniquely qualifies LC-128K on every backend---faster prefill under the same SLO.
But on the cached-RAG workload the ranking flips: the best MI355X config (SGLang,
$99.6$) beats both B200 vLLM and B200 SGLang ($\approx64$), and only B200's
CUDA-only TensorRT-LLM surpasses it ($129.9$). The likely cause is memory capacity:
MI355X carries $288$\,GB HBM3E/GPU against B200's $\sim$192\,GB, giving a larger
resident KV/prefix-cache budget---exactly what LC-Cache rewards. The vendor verdict
is therefore per-workload: B200 for prefill-heavy and ultra-long context, MI355X
competitive-to-better for cached RAG, with the CUDA-only stack (TensorRT-LLM) the
LC-Cache leader where it is available.

\paragraph{Caveats.}
Several capacity points sit at the first ladder rung (e.g.\ LC-32K on the 235B
qualifies only at concurrency~1), so absolute goodput there is small and
prefill-dominated; we report run-to-run variance for boundary points. The 35B
precision comparison is at $\text{TP}{=}1$ (the 35B's multimodal vision-encoder MLP
width, $4304$, is not 16-divisible after an 8-way shard, which dynamic FP8 rejects),
so it should be read against the 35B TP1 BF16 row.
B200 power telemetry was not collected, so its energy columns are unavailable, and
the cross-vendor comparison assumes each platform at its best-performing
configuration (eager on AMD, since graphs regressed there). Finally, the 35B
LC-Cache guardrail failure is treated as a reported result, not grounds for
discarding the runs.
